\documentclass[
     english, % thesis will be in english
    % signatureplaces, % adds signature places in the title page
    monochrome, % makes all link colors black
]{VUMIFSEMasterThesis}
\usepackage{float}
\usepackage{wrapfig2}
\usepackage{hyperref}
\usepackage{algorithmicx}
\usepackage{algorithm}
\usepackage{algpseudocode}
\usepackage{amsfonts}
\usepackage{amsmath}
\usepackage{bm}
\usepackage{caption}
\usepackage{color}
\usepackage{graphicx}
\usepackage{listings}
\usepackage{subcaption}
\usepackage{biblatex}

% Title page
\university{Vilnius University}
\faculty{Faculty of Mathematics and Informatics}
\department{Software Engineering Master's Study Program}
\papertype{Master's Thesis Research Plan}
\title{}
\titleineng{End-to-End Neural Text-to-Speech Synthesis for the Lithuanian Language}
\author{Vytautas Lėveris}
% \secondauthor{Vardonis Pavardonis}   % Add second author
\supervisor{Dr. Gražina Korvel}
\reviewer{...} % Missing
\date{Vilnius – \the\year}

\bibliography{bibliography}

\begin{document}
\maketitle

%% Acknowledgements of individuals and/or institutions.
% \sectionnonumnocontent{}
% \vspace{7cm}
% \begin{center}
%     Padėkos asmenims ir/ar organizacijoms
% \end{center}

%\begin{lithuanian}
%\sectionnonumnocontent{Santrauka}
%Santraukose lietuvių ir anglų kalbomis glaustai aprašomas darbo turinys:
%pristatoma nagrinėta problema ir padarytos išvados. Santraukų gale nurodomi
%darbo raktiniai žodžiai. Santraukos lietuvių ir anglų kalbomis rašomos
%atskiruose puslapiuose. Kiekvienos jų apimtis -- ne daugiau kaip 0,5 puslapio.
%Automatiškai naudojamos lietuviškos kabutės: \enquote{tekstas}.

% Specify up to 5 most important topic keywords.
% A keyword can have multiple words.
%\raktiniaizodziai{raktinis žodis 1, raktinis žodis 2, raktinis žodis 3, raktinis žodis 4, raktinis žodis 5}
%\end{lithuanian}

%\begin{english}
%\sectionnonumnocontent{Summary}
%English summary. English quotes are used automatically: \enquote{tekstas}.

%\keywords{keyword 1, keyword 2, keyword 3, keyword 4, keyword 5}
%\end{english}

%\tableofcontents

\sectionnonumnocontent{Master's Thesis Research Plan}
%Įvade aprašomi darbo tikslai ir uždaviniai, nurodomas temos aktualumas, aptariamos
%teorinės darbo prielaidos bei metodika, apibrėžiamas tiriamasis objektas, apibūdinami
%su tema susiję literatūros ar kitokie šaltiniai, temos analizės tvarka, darbo atlikimo
%aplinkybės, pateikiama žinių apie naudojamus instrumentus (programas ir kt.). Darbo
%įvadas neturi būti dėstymo santrauka. Įvado apimtis 3--4 puslapiai.

\textbf{Introduction}.

\textbf{Background}.

\textbf{Research Object}.

\textbf{Research Goal}.

\textbf{Objectives}:
\begin{enumerate}
\item .
\item .
\item .
\item .
\item .
\end{enumerate}

\textbf{Research Hypothesis}.

\textbf{Research Methods}:
\begin{itemize}
\item ;
\item ;
\item ;
\item .
\end{itemize}

\textbf{Expected Results}:
\begin{enumerate}
\item .
\item .
\end{enumerate}




%\sectionnonum{Rezultatai}
%Rezultatų skyriuje išdėstomi pagrindiniai darbo rezultatai: kažkas išanalizuota,
%kažkas sukurta, kažkas įdiegta. Tarpinių žingsnių išdavos skirtos užtikrinti galutinio
%rezultato kokybę neturi būti pateikiami šiame skyriuje. Kalbant informatikos termi-
%nais, šiame skyriuje pateikiama darbo išvestis, kuri gali būti įvestimi kituose panašios
%tematikos darbuose. Rezultatai pateikiami sunumeruotų (gali būti hierarchiniai) sąrašų
%pavidalu. Darbo rezultatai turi atitikti darbo tikslą.

%\sectionnonum{Išvados}
%\begin{enumerate}[labelindent=0pt]
%    \item Išvadų skyriuje daromi nagrinėtų problemų sprendimo metodų palyginimai, siūlomos
%rekomendacijos, akcentuojamos naujovės.
%    \item Išvados pateikiamos sunumeruoto (gali būti hierarchinis) sąrašo pavidalu.
%    \item Darbo išvados turi atitikti darbo tikslą.
%\end{enumerate}

\printbibliography[heading=bibintoc]  % Šaltinių sąraše nurodoma panaudota
% literatūra, kitokie šaltiniai. Abėcėlės tvarka išdėstomi darbe panaudotų
% (cituotų, perfrazuotų ar bent paminėtų) mokslo leidinių, kitokių publikacijų
% bibliografiniai aprašai. Šaltinių sąrašas spausdinamas iš naujo puslapio.
% Aprašai pateikiami netransliteruoti. Šaltinių sąraše negali būti tokių
% šaltinių, kurie nebuvo paminėti tekste (LaTeX tai sutvarko automatiškai).
% Šaltinių sąraše rekomenduojame necituoti savo kursinio darbo, nes tai nėra
% oficialus literatūros šaltinis. Jei tokių nuorodų reikia, pateikti jas tekste.

% \sectionnonum{Sąvokų apibrėžimai}
%\sectionnonum{Santrumpos}
%Sąvokų apibrėžimai ir santrumpų sąrašas sudaromas tada, kai darbo tekste
%vartojami specialūs paaiškinimo reikalaujantys terminai ir rečiau sutinkamos
%santrumpos.

% Appendices
% Prieduose gali būti pateikiama pagalbinė, ypač darbo autoriaus savarankiškai
% parengta, medžiaga. Savarankiški priedai gali būti pateikiami ir
% kompaktiniame diske. Priedai taip pat numeruojami ir vadinami. Darbo tekstas
% su priedais susiejamas nuorodomis.
%\appendix{Neuroninio tinklo struktūra}

%\begin{figure}[H]
%    \centering
%    \includegraphics[scale=0.5]{img/MLP}
%    \caption{Paveikslėlio pavyzdys}
%    \label{img:mlp}
%\end{figure}


%\appendix{Eksperimentinio palyginimo rezultatai}

% tablesgenerator.com - converts calculators (e.g. excel) tables to LaTeX
%\begin{table}[H]\footnotesize
%  \centering
%  \caption{Lentelės pavyzdys}
%  {\begin{tabular}{|l|c|c|} \hline
%    Algoritmas & $\bar{x}$ & $\sigma^{2}$ \\
%    \hline
%    Algoritmas A  & 1.6335    & 0.5584       \\
%    Algoritmas B  & 1.7395    & 0.5647       \\
%    \hline
%  \end{tabular}}
%  \label{tab:table example}
%\end{table}

\end{document}
